\documentclass[../../main.tex]{subfiles}
\begin{document}

\begin{center}
  {\fontsize{48pt}{58pt}\selectfont \textbf{1979}\normalsize\textbf{年度}}
\end{center}

\begin{table}[htb]
  \centering
  \caption{出題テーマと難易度}
  \begin{tabular}{|c|p{9cm}|c|}
    \hline
    \textbf{問題番号} & \textbf{テーマ} & \textbf{難易度} \\
    \hline
    第1問 & 二つの一次変換による正方形の像 & \\
    \hline
    第2問 & 球と円錐で囲まれた立体の体積 & \\
    \hline
    第3問 & 漸化式で定まる数列と4の倍数 & \\
    \hline
    第4問 & 円内の正三角形と頂点の距離 & \\
    \hline
    第5問 & 三次曲線と直線で囲まれた面積の極限 & \\
    \hline
    第6問 & 垂線を繰り返す点列の収束先 & \\
    \hline
  \end{tabular}
\end{table}

\begin{table}[htb]
  \centering
  \caption{解答}
  \renewcommand{\arraystretch}{1.4}
  \begin{tabular}{|c|c|p{9cm}|}
    \hline
    \textbf{問題番号} & \textbf{小問} & \textbf{答え} \\
    \hline
    第1問 & -- & \\
    \hline
    第2問 & -- & $K$の体積は$\dfrac{2\pi}{3}(11-6\sqrt{3})$ \\
    \hline
    第3問 & -- & $a\equiv-1,0,2\pmod{4}$ \\
    \hline
    第4問 & -- & 最大値は$\dfrac{2}{3}(4-\sqrt{6})$，最小値は$\dfrac{2}{3}$ \\
    \hline
    第5問 & -- & 極限値は$\dfrac{1}{6}$ \\
    \hline
    第6問 & (1) & $P_6=\left(\dfrac{a^2}{(a^2+1)^2},\dfrac{a}{a^2+1}\right)$ \\
    \cline{2-3}
    & (2) & $\left(\dfrac{a^2}{a^4+3a^2+1},\dfrac{a(a^2+1)}{a^4+3a^2+1}\right)$に収束 \\
    \hline
  \end{tabular}
\end{table}

\paragraph{試験形態}
6問．試験時間は要確認．

\end{document}
